\section{}

\paragraph{1309年}\eventanchor{YUAN-1309-REGNAL-YEAR}{YUAN-1309-REGNAL-YEAR}　元以元武宗在位第2年作为诏令、赋役与外交文书的纪年框架。账本所列《元史》〈本纪〉及相关志；《元朝秘史》；《元典章》或《通制条格》相关条为本条来源起点；该条可固定编年位置，具体执行、规模和地方社会影响仍须保留来源差异。

\paragraph{1310年}\eventanchor{YUAN-1310-REGNAL-YEAR}{YUAN-1310-REGNAL-YEAR}　元以元武宗在位第3年作为诏令、赋役与外交文书的纪年框架。账本所列《元史》〈本纪〉及相关志；《元朝秘史》；《元典章》或《通制条格》相关条为本条来源起点；该条可固定编年位置，具体执行、规模和地方社会影响仍须保留来源差异。

\paragraph{1311年}\eventanchor{YUAN-1311-EVENT-183}{YUAN-1311-EVENT-183}　爱育黎拔力八达即位。账本所列《元史》〈本纪〉及相关志；《元朝秘史》；《元典章》或《通制条格》相关条为本条来源起点；该条可固定编年位置，具体执行、规模和地方社会影响仍须保留来源差异。

\paragraph{1311年}\eventanchor{YUAN-1311-REGNAL-YEAR}{YUAN-1311-REGNAL-YEAR}　元以元武宗在位第4年作为诏令、赋役与外交文书的纪年框架。账本所列《元史》〈本纪〉及相关志；《元朝秘史》；《元典章》或《通制条格》相关条为本条来源起点；该条可固定编年位置，具体执行、规模和地方社会影响仍须保留来源差异。

\paragraph{1312年}\eventanchor{YUAN-1312-REGNAL-YEAR}{YUAN-1312-REGNAL-YEAR}　元以元仁宗在位第1年作为诏令、赋役与外交文书的纪年框架。账本所列《元史》〈本纪〉及相关志；《元朝秘史》；《元典章》或《通制条格》相关条为本条来源起点；该条可固定编年位置，具体执行、规模和地方社会影响仍须保留来源差异。

\paragraph{1313年}\eventanchor{YUAN-1313-EVENT-184}{YUAN-1313-EVENT-184}　元廷恢复科举制度。账本所列《元史》〈本纪〉及相关志；《元朝秘史》；《元典章》或《通制条格》相关条为本条来源起点；该条可固定编年位置，具体执行、规模和地方社会影响仍须保留来源差异。

\paragraph{1313年}\eventanchor{YUAN-1313-REGNAL-YEAR}{YUAN-1313-REGNAL-YEAR}　元以元仁宗在位第2年作为诏令、赋役与外交文书的纪年框架。账本所列《元史》〈本纪〉及相关志；《元朝秘史》；《元典章》或《通制条格》相关条为本条来源起点；该条可固定编年位置，具体执行、规模和地方社会影响仍须保留来源差异。

\paragraph{1314年}\eventanchor{YUAN-1314-REGNAL-YEAR}{YUAN-1314-REGNAL-YEAR}　元以元仁宗在位第3年作为诏令、赋役与外交文书的纪年框架。账本所列《元史》〈本纪〉及相关志；《元朝秘史》；《元典章》或《通制条格》相关条为本条来源起点；该条可固定编年位置，具体执行、规模和地方社会影响仍须保留来源差异。

\paragraph{1315年}\eventanchor{YUAN-1315-REGNAL-YEAR}{YUAN-1315-REGNAL-YEAR}　元以元仁宗在位第4年作为诏令、赋役与外交文书的纪年框架。账本所列《元史》〈本纪〉及相关志；《元朝秘史》；《元典章》或《通制条格》相关条为本条来源起点；该条可固定编年位置，具体执行、规模和地方社会影响仍须保留来源差异。

\paragraph{1316年}\eventanchor{YUAN-1316-REGNAL-YEAR}{YUAN-1316-REGNAL-YEAR}　元以元仁宗在位第5年作为诏令、赋役与外交文书的纪年框架。账本所列《元史》〈本纪〉及相关志；《元朝秘史》；《元典章》或《通制条格》相关条为本条来源起点；该条可固定编年位置，具体执行、规模和地方社会影响仍须保留来源差异。

\paragraph{1317年}\eventanchor{YUAN-1317-REGNAL-YEAR}{YUAN-1317-REGNAL-YEAR}　元以元仁宗在位第6年作为诏令、赋役与外交文书的纪年框架。账本所列《元史》〈本纪〉及相关志；《元朝秘史》；《元典章》或《通制条格》相关条为本条来源起点；该条可固定编年位置，具体执行、规模和地方社会影响仍须保留来源差异。

\paragraph{1318年}\eventanchor{YUAN-1318-REGNAL-YEAR}{YUAN-1318-REGNAL-YEAR}　元以元仁宗在位第7年作为诏令、赋役与外交文书的纪年框架。账本所列《元史》〈本纪〉及相关志；《元朝秘史》；《元典章》或《通制条格》相关条为本条来源起点；该条可固定编年位置，具体执行、规模和地方社会影响仍须保留来源差异。

\paragraph{1319年}\eventanchor{YUAN-1319-REGNAL-YEAR}{YUAN-1319-REGNAL-YEAR}　元以元仁宗在位第8年作为诏令、赋役与外交文书的纪年框架。账本所列《元史》〈本纪〉及相关志；《元朝秘史》；《元典章》或《通制条格》相关条为本条来源起点；该条可固定编年位置，具体执行、规模和地方社会影响仍须保留来源差异。

\paragraph{1320年}\eventanchor{YUAN-1320-EVENT-185}{YUAN-1320-EVENT-185}　英宗即位，试图调整宫廷与官僚权力。账本所列《元史》〈本纪〉及相关志；《元朝秘史》；《元典章》或《通制条格》相关条为本条来源起点；该条可固定编年位置，具体执行、规模和地方社会影响仍须保留来源差异。

\paragraph{1320年}\eventanchor{YUAN-1320-REGNAL-YEAR}{YUAN-1320-REGNAL-YEAR}　元以元仁宗在位第9年作为诏令、赋役与外交文书的纪年框架。账本所列《元史》〈本纪〉及相关志；《元朝秘史》；《元典章》或《通制条格》相关条为本条来源起点；该条可固定编年位置，具体执行、规模和地方社会影响仍须保留来源差异。

\paragraph{1321年}\eventanchor{YUAN-1321-REGNAL-YEAR}{YUAN-1321-REGNAL-YEAR}　元以元英宗在位第1年作为诏令、赋役与外交文书的纪年框架。账本所列《元史》〈本纪〉及相关志；《元朝秘史》；《元典章》或《通制条格》相关条为本条来源起点；该条可固定编年位置，具体执行、规模和地方社会影响仍须保留来源差异。

\paragraph{1322年}\eventanchor{YUAN-1322-REGNAL-YEAR}{YUAN-1322-REGNAL-YEAR}　元以元英宗在位第2年作为诏令、赋役与外交文书的纪年框架。账本所列《元史》〈本纪〉及相关志；《元朝秘史》；《元典章》或《通制条格》相关条为本条来源起点；该条可固定编年位置，具体执行、规模和地方社会影响仍须保留来源差异。

\paragraph{1323年}\eventanchor{YUAN-1323-EVENT-186}{YUAN-1323-EVENT-186}　英宗在南坡遇刺，泰定帝即位。账本所列《元史》〈本纪〉及相关志；《元朝秘史》；《元典章》或《通制条格》相关条为本条来源起点；该条可固定编年位置，具体执行、规模和地方社会影响仍须保留来源差异。

\paragraph{1323年}\eventanchor{YUAN-1323-REGNAL-YEAR}{YUAN-1323-REGNAL-YEAR}　元以元英宗在位第3年作为诏令、赋役与外交文书的纪年框架。账本所列《元史》〈本纪〉及相关志；《元朝秘史》；《元典章》或《通制条格》相关条为本条来源起点；该条可固定编年位置，具体执行、规模和地方社会影响仍须保留来源差异。

\paragraph{1324年}\eventanchor{YUAN-1324-REGNAL-YEAR}{YUAN-1324-REGNAL-YEAR}　元以泰定帝在位第1年作为诏令、赋役与外交文书的纪年框架。账本所列《元史》〈本纪〉及相关志；《元朝秘史》；《元典章》或《通制条格》相关条为本条来源起点；该条可固定编年位置，具体执行、规模和地方社会影响仍须保留来源差异。

\paragraph{1325年}\eventanchor{YUAN-1325-REGNAL-YEAR}{YUAN-1325-REGNAL-YEAR}　元以泰定帝在位第2年作为诏令、赋役与外交文书的纪年框架。账本所列《元史》〈本纪〉及相关志；《元朝秘史》；《元典章》或《通制条格》相关条为本条来源起点；该条可固定编年位置，具体执行、规模和地方社会影响仍须保留来源差异。

\paragraph{1326年}\eventanchor{YUAN-1326-REGNAL-YEAR}{YUAN-1326-REGNAL-YEAR}　元以泰定帝在位第3年作为诏令、赋役与外交文书的纪年框架。账本所列《元史》〈本纪〉及相关志；《元朝秘史》；《元典章》或《通制条格》相关条为本条来源起点；该条可固定编年位置，具体执行、规模和地方社会影响仍须保留来源差异。

\paragraph{1327年}\eventanchor{YUAN-1327-REGNAL-YEAR}{YUAN-1327-REGNAL-YEAR}　元以泰定帝在位第4年作为诏令、赋役与外交文书的纪年框架。账本所列《元史》〈本纪〉及相关志；《元朝秘史》；《元典章》或《通制条格》相关条为本条来源起点；该条可固定编年位置，具体执行、规模和地方社会影响仍须保留来源差异。

\paragraph{1328年}\eventanchor{YUAN-1328-EVENT-187}{YUAN-1328-EVENT-187}　元廷爆发两都之战。账本所列《元史》〈本纪〉及相关志；《元朝秘史》；《元典章》或《通制条格》相关条为本条来源起点；该条可固定编年位置，具体执行、规模和地方社会影响仍须保留来源差异。

\paragraph{1328年}\eventanchor{YUAN-1328-REGNAL-YEAR}{YUAN-1328-REGNAL-YEAR}　元以泰定帝在位第5年作为诏令、赋役与外交文书的纪年框架。账本所列《元史》〈本纪〉及相关志；《元朝秘史》；《元典章》或《通制条格》相关条为本条来源起点；该条可固定编年位置，具体执行、规模和地方社会影响仍须保留来源差异。

\paragraph{1329年}\eventanchor{YUAN-1329-REGNAL-YEAR}{YUAN-1329-REGNAL-YEAR}　元以元文宗与元明宗时期在位第1年作为诏令、赋役与外交文书的纪年框架。账本所列《元史》〈本纪〉及相关志；《元朝秘史》；《元典章》或《通制条格》相关条为本条来源起点；该条可固定编年位置，具体执行、规模和地方社会影响仍须保留来源差异。

\paragraph{1330年}\eventanchor{YUAN-1330-REGNAL-YEAR}{YUAN-1330-REGNAL-YEAR}　元以元文宗与元明宗时期在位第2年作为诏令、赋役与外交文书的纪年框架。账本所列《元史》〈本纪〉及相关志；《元朝秘史》；《元典章》或《通制条格》相关条为本条来源起点；该条可固定编年位置，具体执行、规模和地方社会影响仍须保留来源差异。

\paragraph{1331年}\eventanchor{YUAN-1331-REGNAL-YEAR}{YUAN-1331-REGNAL-YEAR}　元以元文宗与元明宗时期在位第3年作为诏令、赋役与外交文书的纪年框架。账本所列《元史》〈本纪〉及相关志；《元朝秘史》；《元典章》或《通制条格》相关条为本条来源起点；该条可固定编年位置，具体执行、规模和地方社会影响仍须保留来源差异。

\paragraph{1332年}\eventanchor{YUAN-1332-REGNAL-YEAR}{YUAN-1332-REGNAL-YEAR}　元以元文宗与元明宗时期在位第4年作为诏令、赋役与外交文书的纪年框架。账本所列《元史》〈本纪〉及相关志；《元朝秘史》；《元典章》或《通制条格》相关条为本条来源起点；该条可固定编年位置，具体执行、规模和地方社会影响仍须保留来源差异。

\paragraph{1333年}\eventanchor{YUAN-1333-EVENT-188}{YUAN-1333-EVENT-188}　妥懽帖睦尔即位。账本所列《元史》〈本纪〉及相关志；《元朝秘史》；《元典章》或《通制条格》相关条为本条来源起点；该条可固定编年位置，具体执行、规模和地方社会影响仍须保留来源差异。

\paragraph{1333年}\eventanchor{YUAN-1333-REGNAL-YEAR}{YUAN-1333-REGNAL-YEAR}　元以元顺帝在位第1年作为诏令、赋役与外交文书的纪年框架。账本所列《元史》〈本纪〉及相关志；《元朝秘史》；《元典章》或《通制条格》相关条为本条来源起点；该条可固定编年位置，具体执行、规模和地方社会影响仍须保留来源差异。

\paragraph{1334年}\eventanchor{YUAN-1334-REGNAL-YEAR}{YUAN-1334-REGNAL-YEAR}　元以元顺帝在位第2年作为诏令、赋役与外交文书的纪年框架。账本所列《元史》〈本纪〉及相关志；《元朝秘史》；《元典章》或《通制条格》相关条为本条来源起点；该条可固定编年位置，具体执行、规模和地方社会影响仍须保留来源差异。

\paragraph{1335年}\eventanchor{YUAN-1335-REGNAL-YEAR}{YUAN-1335-REGNAL-YEAR}　元以元顺帝在位第3年作为诏令、赋役与外交文书的纪年框架。账本所列《元史》〈本纪〉及相关志；《元朝秘史》；《元典章》或《通制条格》相关条为本条来源起点；该条可固定编年位置，具体执行、规模和地方社会影响仍须保留来源差异。

\paragraph{1336年}\eventanchor{YUAN-1336-REGNAL-YEAR}{YUAN-1336-REGNAL-YEAR}　元以元顺帝在位第4年作为诏令、赋役与外交文书的纪年框架。账本所列《元史》〈本纪〉及相关志；《元朝秘史》；《元典章》或《通制条格》相关条为本条来源起点；该条可固定编年位置，具体执行、规模和地方社会影响仍须保留来源差异。

\paragraph{1337年}\eventanchor{YUAN-1337-REGNAL-YEAR}{YUAN-1337-REGNAL-YEAR}　元以元顺帝在位第5年作为诏令、赋役与外交文书的纪年框架。账本所列《元史》〈本纪〉及相关志；《元朝秘史》；《元典章》或《通制条格》相关条为本条来源起点；该条可固定编年位置，具体执行、规模和地方社会影响仍须保留来源差异。

\paragraph{1338年}\eventanchor{YUAN-1338-REGNAL-YEAR}{YUAN-1338-REGNAL-YEAR}　元以元顺帝在位第6年作为诏令、赋役与外交文书的纪年框架。账本所列《元史》〈本纪〉及相关志；《元朝秘史》；《元典章》或《通制条格》相关条为本条来源起点；该条可固定编年位置，具体执行、规模和地方社会影响仍须保留来源差异。

\paragraph{1339年}\eventanchor{YUAN-1339-REGNAL-YEAR}{YUAN-1339-REGNAL-YEAR}　元以元顺帝在位第7年作为诏令、赋役与外交文书的纪年框架。账本所列《元史》〈本纪〉及相关志；《元朝秘史》；《元典章》或《通制条格》相关条为本条来源起点；该条可固定编年位置，具体执行、规模和地方社会影响仍须保留来源差异。

\paragraph{1340年}\eventanchor{YUAN-1340-EVENT-189}{YUAN-1340-EVENT-189}　脱脱掌政后推动修史、钞法和行政整顿。账本所列《元史》〈本纪〉及相关志；《元朝秘史》；《元典章》或《通制条格》相关条为本条来源起点；该条可固定编年位置，具体执行、规模和地方社会影响仍须保留来源差异。

\paragraph{1340年}\eventanchor{YUAN-1340-REGNAL-YEAR}{YUAN-1340-REGNAL-YEAR}　元以元顺帝在位第8年作为诏令、赋役与外交文书的纪年框架。账本所列《元史》〈本纪〉及相关志；《元朝秘史》；《元典章》或《通制条格》相关条为本条来源起点；该条可固定编年位置，具体执行、规模和地方社会影响仍须保留来源差异。

\paragraph{1341年}\eventanchor{YUAN-1341-REGNAL-YEAR}{YUAN-1341-REGNAL-YEAR}　元以元顺帝在位第9年作为诏令、赋役与外交文书的纪年框架。账本所列《元史》〈本纪〉及相关志；《元朝秘史》；《元典章》或《通制条格》相关条为本条来源起点；该条可固定编年位置，具体执行、规模和地方社会影响仍须保留来源差异。

\paragraph{1342年}\eventanchor{YUAN-1342-REGNAL-YEAR}{YUAN-1342-REGNAL-YEAR}　元以元顺帝在位第10年作为诏令、赋役与外交文书的纪年框架。账本所列《元史》〈本纪〉及相关志；《元朝秘史》；《元典章》或《通制条格》相关条为本条来源起点；该条可固定编年位置，具体执行、规模和地方社会影响仍须保留来源差异。

\paragraph{1343年}\eventanchor{YUAN-1343-REGNAL-YEAR}{YUAN-1343-REGNAL-YEAR}　元以元顺帝在位第11年作为诏令、赋役与外交文书的纪年框架。账本所列《元史》〈本纪〉及相关志；《元朝秘史》；《元典章》或《通制条格》相关条为本条来源起点；该条可固定编年位置，具体执行、规模和地方社会影响仍须保留来源差异。

\paragraph{1344年}\eventanchor{YUAN-1344-EVENT-190}{YUAN-1344-EVENT-190}　黄河决溢及相关灾害加深北方社会与财政压力。账本所列《元史》〈本纪〉及相关志；《元朝秘史》；《元典章》或《通制条格》相关条为本条来源起点；该条可固定编年位置，具体执行、规模和地方社会影响仍须保留来源差异。

\paragraph{1344年}\eventanchor{YUAN-1344-REGNAL-YEAR}{YUAN-1344-REGNAL-YEAR}　元以元顺帝在位第12年作为诏令、赋役与外交文书的纪年框架。账本所列《元史》〈本纪〉及相关志；《元朝秘史》；《元典章》或《通制条格》相关条为本条来源起点；该条可固定编年位置，具体执行、规模和地方社会影响仍须保留来源差异。

\paragraph{1345年}\eventanchor{YUAN-1345-REGNAL-YEAR}{YUAN-1345-REGNAL-YEAR}　元以元顺帝在位第13年作为诏令、赋役与外交文书的纪年框架。账本所列《元史》〈本纪〉及相关志；《元朝秘史》；《元典章》或《通制条格》相关条为本条来源起点；该条可固定编年位置，具体执行、规模和地方社会影响仍须保留来源差异。

\paragraph{1346年}\eventanchor{YUAN-1346-REGNAL-YEAR}{YUAN-1346-REGNAL-YEAR}　元以元顺帝在位第14年作为诏令、赋役与外交文书的纪年框架。账本所列《元史》〈本纪〉及相关志；《元朝秘史》；《元典章》或《通制条格》相关条为本条来源起点；该条可固定编年位置，具体执行、规模和地方社会影响仍须保留来源差异。

\paragraph{1347年}\eventanchor{YUAN-1347-REGNAL-YEAR}{YUAN-1347-REGNAL-YEAR}　元以元顺帝在位第15年作为诏令、赋役与外交文书的纪年框架。账本所列《元史》〈本纪〉及相关志；《元朝秘史》；《元典章》或《通制条格》相关条为本条来源起点；该条可固定编年位置，具体执行、规模和地方社会影响仍须保留来源差异。

\paragraph{1348年}\eventanchor{YUAN-1348-EVENT-191}{YUAN-1348-EVENT-191}　方国珍在浙东海域起事。账本所列《元史》〈本纪〉及相关志；《元朝秘史》；《元典章》或《通制条格》相关条为本条来源起点；该条可固定编年位置，具体执行、规模和地方社会影响仍须保留来源差异。

\paragraph{1348年}\eventanchor{YUAN-1348-REGNAL-YEAR}{YUAN-1348-REGNAL-YEAR}　元以元顺帝在位第16年作为诏令、赋役与外交文书的纪年框架。账本所列《元史》〈本纪〉及相关志；《元朝秘史》；《元典章》或《通制条格》相关条为本条来源起点；该条可固定编年位置，具体执行、规模和地方社会影响仍须保留来源差异。

\paragraph{1349年}\eventanchor{YUAN-1349-REGNAL-YEAR}{YUAN-1349-REGNAL-YEAR}　元以元顺帝在位第17年作为诏令、赋役与外交文书的纪年框架。账本所列《元史》〈本纪〉及相关志；《元朝秘史》；《元典章》或《通制条格》相关条为本条来源起点；该条可固定编年位置，具体执行、规模和地方社会影响仍须保留来源差异。

